\documentclass[tikz,border=6pt]{standalone}
\usepackage{mathdiagrams}
% polygon_triangulation_proof — hexagon triangulated from vertex A
% Contract: no white boxes; no prose captions (keep formula only);
% triangle labels must sit in their own triangle, clear of the diagonals.
% A→C has slope 1.9/5.0 = 0.38 — T1 must sit well above that line.
\begin{document}
\begin{tikzpicture}[md]
  \coordinate (A) at (0,0);
  \coordinate (B) at (3.4,0.15);
  \coordinate (C) at (5.0,1.9);
  \coordinate (D) at (4.0,4.0);
  \coordinate (E) at (1.5,4.4);
  \coordinate (F) at (-0.8,2.4);

  \draw[md shape] (A) -- (B) -- (C) -- (D) -- (E) -- (F) -- cycle;
  \draw[md accent, dashed] (A) -- (C);
  \draw[md accent, dashed] (A) -- (D);
  \draw[md accent, dashed] (A) -- (E);

  \node[md point] at (A) {};

  % vertices — plain text, no boxes
  \node[font=\sffamily\small, text=mdInk] at (A) [below left] {$A$};
  \node[font=\sffamily\small, text=mdInk] at (B) [below right] {$B$};
  \node[font=\sffamily\small, text=mdInk] at (C) [right] {$C$};
  \node[font=\sffamily\small, text=mdInk] at (D) [above right] {$D$};
  \node[font=\sffamily\small, text=mdInk] at (E) [above left] {$E$};
  \node[font=\sffamily\small, text=mdInk] at (F) [left] {$F$};

  % triangle labels — inside each triangle, clear of the dashed diagonals
  % T1 = A-B-C; A→C slope 0.38, at x=1.6 the line is y=0.61 → put T1 at y=1.10
  \node[font=\sffamily\small, text=mdAccent] at (1.6,1.10) {$T_1$};
  % T2 = A-C-D; between A→C (y=0.38x) and A→D (y=x)
  \node[font=\sffamily\small, text=mdAccent] at (3.0,1.85) {$T_2$};
  % T3 = A-D-E; above A→D (y=x), left of D→E
  \node[font=\sffamily\small, text=mdAccent] at (2.55,3.35) {$T_3$};
  % T4 = A-E-F; left of A→E (steep)
  \node[font=\sffamily\small, text=mdAccent] at (0.75,2.95) {$T_4$};

  % formula only — no prose prefix
  \node[font=\sffamily\small, text=mdInk] at (2.2,-1.05)
    {$d = n - 3 \quad t = n - 2$};
\end{tikzpicture}
\end{document}
